% !TITLE 虞美人·春花秋月何时了
% !AUTHOR 李煜
% !DYNASTY 五代十国
% !GENRE 词
% !ORDER 4

\begin{poem}
春花秋月何时了\textsuperscript{1}？往事知多少。
小楼昨夜又东风\textsuperscript{2}，故国\textsuperscript{3}不堪\textsuperscript{4}回首月明中。

雕栏玉砌\textsuperscript{5}应犹在，只是朱颜\textsuperscript{6}改。
问君能有几多愁？恰似\textsuperscript{7}一江春水向东流。
\end{poem}

\begin{annotations}
\notetext{1}{了：结束、终结。春花秋月本是美好的事物，李煜却问它们什么时候结束——因为美景只会让他想起已经永远失去的光景。欢愉的事物成了痛苦的触发点。}
\notetext{2}{东风：春天的风。昨夜又是春风吹来——年复一年，春风吹拂，但故国已经不再了。}
\notetext{3}{故国：已经灭亡的南唐。被俘后身在汴京，不敢朝月光的方向望——因为月明的方向，就是故国的方向。}
\notetext{4}{不堪：承受不了。无法承受在月光中回想故国的痛苦。}
\notetext{5}{雕栏玉砌：雕刻的栏杆、玉石砌的台阶。指南京南唐皇宫中的建筑。建筑物应该还在，但主人已经换了。}
\notetext{6}{朱颜：红润的面容。也代指故国的宫女们，以及自己的青春和帝王生涯——一切都已改变。}
\notetext{7}{恰似：正好像。愁思像一江春水，滚滚东流，永不枯竭。把抽象的愁比作可见的江水，这个比喻让千百年来的读者都能“看见”他的痛苦。}
\end{annotations}

\begin{appreciation}
这是李煜最著名的词，也是整个词史上传唱最广的词之一。据说写完这首词不久，他就被宋朝的皇帝赐了毒酒——一首词，要了一个人的命，在中国文学史上，这是头一回。

“春花秋月何时了？”开篇就是一句问。春花、秋月，人间最美的两样东西，他问它们什么时候才结束。因为美好的事物只会提醒他，那些美好已经永远失去了。常人赏花赏月，李煜怕花怕月。“往事知多少”，一问之下，往事全涌上来。

“小楼昨夜又东风”，一个“又”字，是说东风已经不是第一年吹来了。春天年年回来，故国却不再回来。“故国不堪回首月明中”——月亮的方向，就是故乡的方向，他不敢朝那个方向望。“雕栏玉砌应犹在”，皇宫里的栏杆台阶应该还在，只是主人换了；“只是朱颜改”，宫女的容颜改了，他自己的青春和帝王生涯也改了。屈原在《离骚》里说“恐美人之迟暮”，怕美好的人和事老去——李煜这里写的，是眼睁睁看着这一切发生，连怕的资格都没有了。

最后是自问自答：“问君能有几多愁？恰似一江春水向东流。”愁有多少？像一江春水向东流。这个比喻好在哪里？好在春水：春天的江水涨起来，浩荡无边，拦不住，而且向东流——东边，正是故国的方向。连愁的方向，都是回家的方向。王国维在《人间词话》里评他：词到李后主，眼界始大，感慨遂深，把歌女唱的词，变成了文人写自己的心。这一句“一江春水”，就是那个转变的见证。

这首词写完不久，李煜就死了，死在汴京，只有四十一岁。一江春水还在流，他没能再看一眼故国的花。要我说，春花秋月还是要好好看的，它们没有错，错的是把它们看丢了的人。今年春天，记得带妈妈去公园看花。
\end{appreciation}
